\documentclass[11pt]{article}

\usepackage[margin=1in]{geometry}
\usepackage{amsmath,amssymb,amsthm,mathtools}
\usepackage{hyperref}

\newcommand{\R}{\mathbb{R}}
\newcommand{\Sym}{\operatorname{Sym}}

\newtheorem{theorem}{Theorem}
\newtheorem{proposition}[theorem]{Proposition}
\newtheorem{corollary}[theorem]{Corollary}
\theoremstyle{remark}
\newtheorem{remark}[theorem]{Remark}
\newtheorem{example}[theorem]{Example}

\title{Do Boundary Fast Weights Persist or Collapse?}
\author{Research note generated in sandbox}
\date{2026-04-05}

\begin{document}
\maketitle

\section{Question}

Let $q_{\alpha,t}\in\R^p$ denote the member-specific boundary fast state in the mean-field limit, with dynamics
\begin{equation}
\dot q_{\alpha,t}=\mathcal C(\rho_t,q_{\alpha,t}),
\qquad
\alpha=1,\dots,K.
\label{eq:q-dynamics}
\end{equation}
If the members are initialized with different boundary fast weights, must these differences disappear during training?

\medskip

Answer:

\begin{quote}
No. Different boundary initializations disappear only in a contractive regime. In general they may decay, persist, or grow.
\end{quote}

\section{A Clean Necessary Principle}

Fix two members $\alpha,\beta$ and define
\[
\Delta_t:=q_{\alpha,t}-q_{\beta,t}.
\]
Then
\[
\dot\Delta_t
=
\mathcal C(\rho_t,q_{\alpha,t})-\mathcal C(\rho_t,q_{\beta,t}).
\]

\begin{theorem}[Uniform contraction implies disappearance]
\label{thm:uniform-contraction}
Assume there exists $\gamma>0$ such that for every $\rho\in\mathcal X$ and every $q,q'\in\R^p$,
\begin{equation}
\langle q-q',\,\mathcal C(\rho,q)-\mathcal C(\rho,q')\rangle
\le
-\gamma \|q-q'\|^2.
\label{eq:onesided}
\end{equation}
Then
\[
\|q_{\alpha,t}-q_{\beta,t}\|
\le
e^{-\gamma t}\|q_{\alpha,0}-q_{\beta,0}\|
\qquad
\forall t\ge 0.
\]
In particular, all boundary-state differences vanish exponentially.
\end{theorem}

\begin{proof}
Differentiate $\frac12\|\Delta_t\|^2$ and use \eqref{eq:onesided}:
\[
\frac{d}{dt}\frac12\|\Delta_t\|^2
=
\langle \Delta_t,\dot\Delta_t\rangle
\le
-\gamma \|\Delta_t\|^2.
\]
Gronwall gives the result.
\end{proof}

\begin{remark}
So if one can prove a strong one-sided Lipschitz condition for the boundary-state drift, then different boundary fast weights \emph{do} wash out. But this is a genuine assumption, not a universal fact.
\end{remark}

\section{Why Differences Need Not Disappear}

\begin{example}[Exact persistence]
\label{ex:persistence}
Take
\[
\mathcal C(\rho,q)=b(\rho),
\]
independent of $q$. Then every member obeys the same translated dynamics
\[
\dot q_{\alpha,t}=b(\rho_t),
\]
so
\[
\frac{d}{dt}(q_{\alpha,t}-q_{\beta,t})=0.
\]
Hence
\[
q_{\alpha,t}-q_{\beta,t}=q_{\alpha,0}-q_{\beta,0}
\qquad
\forall t.
\]
Initial boundary differences persist exactly forever.
\end{example}

\begin{example}[Growth]
\label{ex:growth}
Take
\[
\mathcal C(\rho,q)=Aq
\]
with a matrix $A$ having an eigenvalue of positive real part. Then differences evolve by
\[
\dot\Delta_t=A\Delta_t,
\]
and components in unstable eigendirections grow exponentially.
\end{example}

\begin{example}[Multiple attractors]
\label{ex:multistable}
In one dimension, let
\[
\mathcal C(\rho,q)=q-q^3.
\]
Then the dynamics has stable equilibria at $\pm 1$. If one member starts positive and another negative, they converge to different terminal states rather than collapsing together.
\end{example}

\begin{corollary}[No universal disappearance theorem]
\label{cor:no-universal}
Without an additional contraction assumption on $\mathcal C$, there is no theorem stating that different boundary initializations must vanish in the mean-field limit.
\end{corollary}

\section{Local Stationary-Law Test}

Suppose the collapsed system converges to a stationary point $(\rho_\star,q_\star)$ with
\[
\mathcal C(\rho_\star,q_\star)=0.
\]
Define
\[
A_\star:=\partial_q \mathcal C(\rho_\star,q_\star).
\]

\begin{proposition}[Linear stationary-law criterion]
\label{prop:linear-criterion}
Near $(\rho_\star,q_\star)$, the fate of small boundary differences is determined at first order by
\[
\dot\Delta_t=A_\star \Delta_t.
\]
Hence:
\begin{itemize}
\item if $\Sym(A_\star)\preceq -\gamma I$ for some $\gamma>0$, then small differences decay exponentially;
\item if $A_\star$ has an eigenvalue with positive real part, then some small differences grow;
\item if $A_\star$ has neutral directions, then first-order theory alone does not force disappearance.
\end{itemize}
\end{proposition}

\begin{proof}
Linearize $\mathcal C(\rho_\star,q)$ in $q$ around $q_\star$:
\[
\mathcal C(\rho_\star,q_\star+\Delta)
=
A_\star \Delta + o(\|\Delta\|).
\]
This gives the stated linearized equation. The three bullets are standard consequences of linear ODE theory.
\end{proof}

\begin{remark}
This criterion is the precise answer to the user's question in local mean-field terms. Boundary differences disappear only if the terminal boundary dynamics is stable in the member-specific directions.
\end{remark}

\section{What the Current Theory Can and Cannot Say}

\begin{corollary}[Safe statement]
\label{cor:safe}
The current mean-field theory supports the following claim:

\begin{quote}
Different boundary fast weights may persist, decay, or grow. To know which occurs, one must analyze the boundary-state drift $\mathcal C$ or its stationary Jacobian $A_\star$.
\end{quote}
\end{corollary}

\begin{remark}
So if your intuition is ``boundary fast weights might not disappear'', that intuition is mathematically valid. They disappear only in a sufficiently contractive regime. Otherwise, they can survive all the way to the terminal law and may even be the correct place to encode finite-$K$ diversity.
\end{remark}

\section{Practical Consequence}

There are two qualitatively different regimes:
\begin{itemize}
\item \textbf{Contractive boundary regime.}
  Initialization-only diversity in the boundary variables is transient. One needs ongoing centered forcing, explicit constraints, or redesign at the terminal law.
\item \textbf{Persistent boundary regime.}
  Boundary differences can survive to the end. In that case, stationary-law simplex/tetrahedral design is meaningful as an actual terminal predictor design, not just a transient initialization trick.
\end{itemize}

\end{document}
